\paragraph{SHAP (SHapley Additive exPlanations)}
SHAP unifies several additive feature-attribution methods under a single game-theoretic framework,
which treats each input feature like a player in a cooperative game and fairly shares out the credit
for a prediction among them, assigning each feature a Shapley value that quantifies its marginal
contribution to an individual prediction \cite{Lundberg2017}. It satisfies desirable properties such as
local accuracy and consistency, and has efficient model-specific variants such as TreeSHAP and
DeepSHAP. However, SHAP produces continuous, instance-level attributions, that is, a numeric importance
score for each feature on a single prediction rather than a discrete yes/no rule: it reports which genes
mattered for one prediction, not a generalisable decision rule that can be audited globally or expressed
as a discrete gene set for cross-referencing against KEGG pathways. Most critically for this project,
attribution-based methods of this kind have been shown to produce visually plausible explanations even
for models with randomised weights \cite{Adebayo2018}, so a SHAP explanation of an overfitted model can
look indistinguishable from that of a well-generalised one; SHAP therefore cannot detect the Clever Hans
behaviour this pipeline targets. Adversarial work, meaning deliberately crafted inputs or training
procedures designed to fool an explanation method, has further shown that SHAP can be manipulated to
conceal biased model behaviour \cite{Slack2020}. SHAP is predominantly deployed in Blue-XAI practice,
prioritising end-user justification over model interrogation.

\paragraph{LIME (Local Interpretable Model-Agnostic Explanations)}
LIME explains an individual prediction by sampling perturbed instances around it and fitting a sparse,
proximity-weighted linear surrogate (a simple linear model fitted only in the local region around the
prediction), yielding local feature importances \cite{Ribeiro2016}. Because the explanation is a local
linear approximation, it is inherently per-instance rather than global, and its faithfulness degrades
where the underlying decision surface is highly non-linear. The method is also unstable: small changes
to the perturbation neighbourhood can produce qualitatively different explanations
\cite{AlvarezMelis2018}, and in high-dimensional genomic feature spaces meaningful local neighbourhoods
become sparse, undermining the sampling strategy. Like SHAP, LIME is vulnerable to adversarial
manipulation that hides discriminatory behaviour \cite{Slack2020}, and its continuous attributions are
not directly verifiable against discrete biological pathways. LIME is thus predominantly a Blue-XAI tool
for per-prediction justification.

\paragraph{Layer-wise Relevance Propagation (LRP)}
LRP explains a prediction by back-propagating the model output and redistributing it backwards as a
relevance signal, a share of the prediction assigned to each input, onto the input features, producing a
heatmap, a colour-coded map of which inputs mattered most \cite{Bach2015,Samek2017}. While effective for
visualising salient inputs, it yields continuous relevance maps rather than auditable symbolic rules and
remains local to individual predictions. Most importantly for this project, LRP is the very instrument
through which the Clever Hans problem has been demonstrated: Pahde et al.~\cite{Pahde2023} apply LRP to
overfitted models and show it attributes high relevance to biologically incoherent or artefactual
features the model has latched onto. LRP therefore functions here as a witness to the fragility being
diagnosed rather than as a competing rule-extraction method.

\paragraph{Integrated Gradients}
Integrated Gradients attributes a prediction to its input features by accumulating gradients along a
straight-line path from a baseline input to the actual input, satisfying the completeness and
sensitivity axioms, formal guarantees that the attributions sum to the prediction and that any feature
which changes the output receives non-zero importance, that naive gradient saliency (reading the model's
raw input gradients directly) violates \cite{Sundararajan2017}. Although axiomatically stronger than
these vanilla gradient methods, it shares their core limitations: the output is a continuous,
instance-level attribution rather than a global rule, the result depends on the often-arbitrary choice
of baseline, and the path integral, the accumulation of gradients along that path, is sensitive to the
smoothness of the model's decision boundary. This last property is precisely the sharp-loss-landscape
pathology that the SAM intervention in this project is designed to correct, meaning Integrated Gradients
is most fragile exactly where model health is poorest.

\paragraph{Anchors}
Anchors is the closest state-of-the-art method to the approach taken here: rather than continuous
attributions, it produces high-precision if-then rules, where an ``anchor'' is a set of feature
conditions that locally guarantees the prediction to a user-specified precision \cite{Ribeiro2018}. This
rule-like, human-readable output is more aligned with clinical auditability than attribution heatmaps.
However, anchors are inherently local: one rule is generated per instance, with no global rule set
describing the model's overall logic, and the precision guarantee holds only in the neighbourhood of
each instance, so it need not generalise across the heterogeneous cancer subtypes in GDSC. Anchors also
optimises only for precision and coverage; it does not expose fidelity and comprehensibility as explicit
competing objectives, cannot produce a Pareto front of trade-off solutions, and therefore cannot feed
the PROMETHEE II multi-criteria ranking that structures this project's evaluation.
\\
The limitations of these SOTA methods motivate the use of evolutionary rule extraction in this project.
Unlike SHAP, LIME, LRP, and Integrated Gradients, which produce continuous feature-attribution scores
that cannot be directly verified against biological pathways in a discrete, auditable form, evolutionary
algorithms extract global symbolic if-then rules that can be explicitly evaluated for fidelity,
comprehensibility, and biological validity \cite{Yedjour2024}. Unlike Anchors, the GA and NSGA-II
pipelines optimise rule quality as explicit objectives across the full GDSC dataset rather than locally
per instance, and their outputs are directly compatible with the PROMETHEE II multi-criteria ranking
framework \cite{Yedjour2024}. Furthermore, and most critically, none of the SOTA methods described above
can detect Clever Hans behaviour: Adebayo et al.~\cite{Adebayo2018} demonstrated that popular attribution
methods including SHAP-style approaches produce visually plausible explanations even for models that have
memorised noise, because their outputs are largely independent of model parameters. This project's Red
XAI pipeline addresses this directly by auditing model health before any explanation is extracted
\cite{Biecek2024}. Table~\ref{tab:sota_comparison} provides a structured comparison.

\begin{table}[H]
\centering
\small
\caption{Structured comparison of state-of-the-art post-hoc XAI methods against the evolutionary
rule-extraction approach adopted in this project.}
\label{tab:sota_comparison}
\begin{tabular}{p{1.9cm}p{1.7cm}p{1.0cm}p{1.5cm}p{1.3cm}p{1.6cm}p{1.1cm}}
\toprule
\textbf{Method} & \textbf{Output} & \textbf{Scope} & \textbf{Auditable Rules} & \textbf{Fidelity Opt.} & \textbf{Bio. Verifiable} & \textbf{Red XAI} \\
\midrule
SHAP & Continuous scores & Local & No & No & Indirect & No \\
LIME & Continuous scores & Local & No & No & Indirect & No \\
LRP & Heatmap & Local & No & No & Indirect & No \\
Integrated Gradients & Continuous scores & Local & No & No & Indirect & No \\
Anchors & Rule (local) & Local & Partial & No & Indirect & Partial \\
\textbf{GA + NSGA-II (Ours)} & \textbf{Symbolic rules} & \textbf{Global} & \textbf{Yes} & \textbf{Yes} & \textbf{Direct (GSEA)} & \textbf{Yes} \\
\bottomrule
\end{tabular}
\end{table}

\noindent
Beyond these post-hoc methods, a distinct but related family known as explanation-guided training
modifies the training objective itself rather than explaining a fixed model after the fact. Examples
include Right for the Right Reasons \cite{Ross2017}, contextual-decomposition penalisation
\cite{Rieger2020}, adversarial training on explanations \cite{Tang2022}, and robust-explanation
objectives such as R2ET \cite{Chen2022}, REGEX \cite{Li2023}, and the Clever-Hans-specific pruning
method EGEM \cite{Linhardt2024}. These methods place the explanation directly into the loss function, in
some cases requiring human-annotated ground-truth explanations or expensive second-order gradients. They
are the nearest methodological neighbour to the approach proposed here and are revisited in the critical
analysis (Section~\ref{ssub:critical-analysis}), where their relationship to the explanation-agnostic,
geometry-only intervention used in this project is examined.